%==================================================================
% BAB 2 — TINJAUAN PUSTAKA
%==================================================================

\chapter{TINJAUAN PUSTAKA}

\section{Tinjauan Pustaka}

% CATATAN PENTING
% Gunakan MATRIKS SINTESIS, bukan tabel per-penelitian dengan kolom
% "perbedaan dengan penelitian ini". Pada matriks, celah penelitian
% TERLIHAT sebagai pola (kolom yang kosong di semua baris), sehingga
% Anda cukup menuliskan apa yang tampak — tidak perlu mengarang kebaruan.
%
% Kolom dipilih berdasarkan relevansinya dengan rumusan masalah Anda.
% Empat sampai enam penelitian yang benar-benar dibedah jauh lebih baik
% daripada lima belas yang sekadar didaftar.

Bagian ini menyintesiskan penelitian-penelitian terdahulu yang relevan dengan penanganan latensi dialog NPC, antarmuka diegetik, dan pengalaman pengguna (\textit{user experience}). Sintesis disajikan pada \cref{tab:sintesis} untuk memperlihatkan posisi dan celah penelitian.

\begin{table}[H]
    \centering
    \caption{Matriks sintesis penelitian terdahulu}
    \label{tab:sintesis}
    \footnotesize
    \setlength{\tabcolsep}{4pt}
    \begin{tabular}{|p{2.2cm}|p{2.5cm}|p{2.5cm}|p{2.2cm}|p{2.5cm}|p{2.3cm}|}
        \hline
        \textbf{Penelitian} & \textbf{Konteks/Objek \textit{Game}} & \textbf{Sumber/Beban Latensi} & \textbf{Bentuk Antarmuka Pemrosesan} & \textbf{Mekanisme \textit{Latency Masking}} & \textbf{Metrik Evaluasi} \\ \hline

        \citet{siahaan2026retrieval} & Dialog NPC \& generasi \textit{quest} pada RPG C++ & Pemrosesan inferensi model lokal (23,5--220,5 detik) & \textit{Non-Diegetic} (\textit{Console Output}) & Tidak ada penyamaran latensi & Hanya evaluasi akurasi RAG \\ \hline

        \citet{wei2025mind} & NPC pembangun pada \textit{game survival} 3D \textit{Echoshell} & Pemrosesan kognitif LLM & \textit{Diegetic} \& \textit{Non-Diegetic} & Pemisahan \textit{thread Mind \& Body} & Tidak mengukur \textit{perceived latency} \\ \hline

        \citet{elfleet2024investigating} & Agen percakapan bodi di \textit{Virtual Reality} (VR) & Latensi transkripsi ASR dan NLU LLM & \textit{Non-Diegetic} (\textit{Spinner UI}) & \textit{Filled pauses} \& animasi nonverbal & \textit{Perceived Latency} \& \textit{Immersion} \\ \hline
        
        \citet{csepregi2021effect} & Dialog NPC RPG berbasis LLM \textit{Context-Aware} & Waktu tunggu \textit{generation response} & \textit{Diegetic} (\textit{Embed-based animation} \& Audio) & Memicu animasi dari frasa yang di-\textit{embed} & \textit{Player Engagement} pada \textit{Context Aware NPC} \\ \hline

        \citet{andersson2024effects} & Pengaruh UI diegetik vs \textit{non-diegetic} pada \textit{game demo} & Pemrosesan sistem \textit{game} statis & \textit{Diegetic} UI & Hanya untuk menu/inventaris & Mengukur \textit{Immersion \& Usability} \\ \hline

        \textbf{Penelitian Ini} & NPC Pustakawan \textit{Game} 2D Eksplorasi Kampus & Jeda Respons NPC (\textit{Diegetic Masking} vs \textit{Non-Diegetic Baseline}) & \textit{Diegetic} & Animasi Kontekstual (\textit{Asynchronous}) & Pengujian \textit{Perceived Latency, Immersion}, \&  retensi \\ \hline
    \end{tabular}
\end{table}

% Paragraf berikut MEMBACA POLA pada matriks, lalu menyatakan posisi
% penelitian. Perhatikan bahwa celah dinyatakan berdasarkan apa yang
% terlihat pada tabel, bukan diklaim begitu saja.
Berdasarkan \cref{tab:sintesis}, terlihat pola celah penelitian (\textit{research gap}) yang mendasari posisi penelitian ini. Pertama, pola celah efektivitas penyamaran latensi pada NPC 2D. Penelitian \citet{siahaan2026retrieval} mengidentifikasi bahwa pemrosesan dialog NPC membutuhkan durasi waktu komputasi, hal ini memicu masalah agen yang membeku (\textit{frozen agent}). \citet{elfleet2024investigating} menggunakan \textit{spinner} UI pada kepala NPC dan animasi \textit{idle} (nonverbal) untuk mengurangi latensi di VR, sementara \citet{csepregi2021effect} men- \textit{trigger} animasi di RPG. Namun, kedua studi tersebut berfokus pada lingkungan 3D atau VR (ECA), sehingga belum ada adaptasi mekanisme \textit{masking} berbasis animasi kontekstual pada ruang lingkup \textit{game engine} 2D.

Kedua, pola batasan pemanfaatan antarmuka diegetik. \citet{andersson2024effects} membuktikan bahwa \textit{Diegetic} UI efektif meningkatkan \textit{player immersion}. Namun, penerapannya sejauh ini hanya terbatas pada elemen visual statis seperti inventaris, tanpa pernah dimanfaatkan secara fungsional sebagai elemen rekayasa (\textit{latency masking}) saat waktu komputasi terjadi di latar belakang.

Ketiga, pola keterpisahan pengukuran \textit{perceived latency} dengan eksekusi \textit{Mind \& Body}. Walaupun \citet{wei2025mind} berhasil memisahkan proses kognitif dan eksekusi fisik NPC (\textit{Mind and Body}), penelitian tersebut tidak mengukur dampak persepsi waktu tunggu (\textit{perceived latency}) secara terisolasi. Di sisi lain, \citet{elfleet2024investigating} dan \citet{csepregi2021effect} membuktikan secara kuantitatif melalui \textit{playtesting} bahwa umpan balik multimodal meningkatkan \textit{immersion} dan \textit{player engagement}.

Penelitian ini mengisi celah tersebut dengan mengimplementasikan animasi kontekstual NPC secara \textit{asynchronous} sebagai bentuk \textit{Diegetic Latency Masking} pada \textit{game} 2D, serta menguji dampaknya terhadap \textit{perceived latency}, \textit{player immersion}, dan retensi pemain dibandingkan indikator \textit{non-diegetic} konvensional.

\section{Landasan Teori}
% Susun subbab sesuai konsep utama penelitian Anda. Tiap subbab memuat
% definisi, unsur pokok, dan rujukannya. Hindari menyalin definisi
% panjang; ringkas dengan bahasa sendiri lalu sitasi sumbernya.

\subsection{\textit{Diegetic User Interface} (\textit{Diegetic} UI)}
Antarmuka diegetik (\textit{Diegetic User Interface}) adalah elemen antarmuka pengguna yang terdapat dan terintegrasi secara fisik serta geometris di dalam dunia fiksi cerita permainan, sehingga dapat dilihat dan diinteraksi oleh karakter di dalam \textit{game} maupun oleh pemain \citep{andersson2024effects}. Berbeda dengan antarmuka \textit{non-diegetic} (seperti \textit{overlay} HUD, \textit{loading bar}, atau teks indikator pemuatan), antarmuka diegetik mempertahankan kehadiran spasial (\textit{spatial presence}) dan imersi naratif (\textit{player immersion}) karena tidak menarik perhatian pemain keluar dari dunia fiksi \citep{andersson2024effects}.

\subsection{\textit{Latency Masking} dan \textit{Perceived Latency} dalam HCI}
Dalam konteks \textit{Human-Computer Interaction} (HCI), latensi dapat dibedakan menjadi \textit{objective latency} (durasi waktu tunggu riil dalam detik) dan \textit{perceived latency} (durasi waktu tunggu yang dirasakan secara subjektif oleh pengguna) \citep{elfleet2024investigating}. \textit{Latency Masking} adalah teknik rekayasa antarmuka untuk memanipulasi perhatian pengguna (\textit{attentional bias}) selama jeda komputasi berlangsung. Dengan memberikan umpan balik visual yang bermakna seperti pergerakan \textit{nonverbal} atau animasi kontekstual di dalam dunia permainan, \textit{perceived latency} dapat dikurangi secara signifikan \citep{elfleet2024investigating}.

\subsection{\textit{Asynchronous State} dan \textit{Execution} dalam \textit{Game Engine}}
Dalam arsitektur \textit{game engine}, eksekusi logika permainan berjalan berbasis \textit{frame-by-frame}. Untuk mencegah kondisi \textit{frozen agent} saat sistem sedang menunggu pemrosesan respons di latar belakang, \textit{game engine} menerapkan pemisahan \textit{thread} secara \textit{asynchronous} \citep{wei2025mind}. Melalui pola ini, \textit{state machine} animasi fisik NPC tetap mengeksekusi \textit{loop} animasi kontekstual (seperti pergerakan navigasi ke koordinat tertentu) di \textit{main thread}, sementara proses komputasi respons berjalan secara \textit{asynchronous} di \textit{background thread} \citep{wei2025mind}.

\subsection{\textit{Player Immersion} dan \textit{Player Experience} (PX)}
\textit{Player immersion} merujuk pada tingkat keterlibatan mental, emosional, dan persepsi pemain di mana ia merasa terhubung secara mendalam dengan dunia fiksi permainan \citep{andersson2024effects, elfleet2024investigating}. Penggunaan antarmuka \textit{non-diegetic} konvensional saat jeda waktu komputasi sering kali memutus aliran pengalaman (\textit{flow}) dan menurunkan keterlibatan pemain (\textit{player engagement}) \citep{csepregi2021effect}. Penerapan elemen antarmuka maupun dialog yang disesuaikan secara kontekstual dengan narasi terbukti dapat mempertahankan \textit{player immersion} dan meningkatkan kualitas pengalaman bermain (\textit{player experience}) secara keseluruhan \citep{andersson2024effects, csepregi2021effect}.

\subsection{\textit{Game Development Life Cycle} (GDLC)}
\textit{Game Development Life Cycle} (GDLC) adalah metodologi terstruktur yang menjadi panduan standar dalam proses rekayasa perangkat lunak, secara khusus untuk pengembangan \textit{video game} \citep{ramadan2013game}. Siklus hidup ini secara umum mencakup fase inisiasi, praproduksi, produksi, pengujian (\textit{testing}), fase \textit{beta}, hingga rilis \citep{ramadan2013game}. Dalam konteks penelitian ini, pendekatan GDLC diterapkan secara spesifik pada fase praproduksi dan pengujian purwarupa (\textit{prototyping}), di mana rekayasa mekanik \textit{latency masking} diimplementasikan ke dalam purwarupa 2D untuk kemudian divalidasi kelayakannya melalui pengujian pengalaman pengguna (\textit{playtesting}).

\section{Justifikasi Pemilihan Metode}
% SUBBAB INI SERING DILUPAKAN, PADAHAL PALING PENTING.
% Deskriptor KKNI jenjang 6 menuntut kemampuan memilih alternatif solusi
% secara sadar. Di sinilah kemampuan itu dibuktikan.
% Bandingkan minimal dua alternatif, lalu nyatakan alasan pemilihan Anda.

Sesuai dengan standar analisis kualifikasi KKNI Jenjang 6, subbab ini memberikan justifikasi rasional ilmiah terhadap alternatif solusi yang dipilih:

\subsection{Justifikasi Pemilihan Antarmuka Penanganan Latensi}
\begin{itemize}
    \item \textbf{Alternatif A: \textit{Non-Diegetic Indicator} (\textit{Loading Bar} / Teks "\textit{Thinking...}")}. Menampilkan komponen visual antarmuka konvensional di atas layar permainan saat sistem sedang memproses respons. Kelemahan pendekatan ini adalah memutus alur imersi naratif, membuat waktu tunggu terasa lebih lama karena pengguna difokuskan pada proses komputasi, dan secara eksplisit mengingatkan pemain akan keterbatasan teknis sistem \citep{andersson2024effects, elfleet2024investigating}.
    \item \textbf{Alternatif B (Metode Terpilih): \textit{Diegetic Latency Masking} (Animasi Fisik Kontekstual NPC)}. Mengalihkan jeda waktu tunggu komputasi menjadi gerakan fisik kontekstual NPC di dalam dunia 2D. Metode ini dipilih karena mampu menyamarkan \textit{perceived latency} dengan membingkai waktu komputasi sebagai tindakan akting alami karakter di dalam cerita (misal: mencari buku), sehingga menjaga kelancaran \textit{player immersion} dan menstimulasi peningkatan \textit{player engagement} \citep{andersson2024effects, csepregi2021effect, wei2025mind}.
\end{itemize}

\subsection{Justifikasi Pemilihan Metode Pengujian Pengalaman Pengguna}
\begin{itemize}
    \item \textbf{Metode Terpilih: A/B \textit{Testing} dengan \textit{Playtesting} \& \textit{Game Experience Questionnaire} (GEQ)}. Menguji dua kelompok perlakuan (Skenario A: \textit{Non-Diegetic} vs Skenario B: \textit{Diegetic Masking}) dengan durasi jeda waktu komputasi yang diisolasi (\textit{controlled latency}), dilanjutkan dengan pengisian instrumen \textit{Game Experience Questionnaire} (GEQ) \citep{ijsselsteijn2013game}. Penelitian ini berfokus pada \textit{Core Module} GEQ, secara khusus pada dimensi \textit{Sensory and Imaginative Immersion} (mengukur tingkat imersi), \textit{Flow} (mengukur kontinuitas pengalaman), dan \textit{Tension/Annoyance} (mengukur frustrasi akibat latensi), ditambah butir khusus untuk \textit{perceived latency}. Pendekatan ini dipilih karena GEQ memberikan instrumen empiris kuantitatif yang telah divalidasi khusus untuk membandingkan secara langsung kualitas pengalaman psikologis pemain antar-skenario dalam konteks \textit{video game} \citep{csepregi2021effect, elfleet2024investigating, ijsselsteijn2013game}.
\end{itemize}